\documentclass[10pt,letterpaper]{article}
\usepackage[margin=0.65in]{geometry}
\usepackage{booktabs}
\usepackage{longtable}
\usepackage{graphicx}
\usepackage{hyperref}
\title{Global Phylogenomic Radiation of Isoptera and Sister Wood-Roach Outgroups (Blattodea: Termitoidae): Native Rust Quartet Coalescent Species-Tree Inference, Endogenous Lignocellulase and Cytochrome P450 Selection, and Cross-Regional vs. Multi-Locus Demographic Skylines}
\author{Bryan P. White \\ CRV Science Labs}
\date{November 10, 2026 \quad|\quad Preprint ID: PREPRINT-26 (v2.0.0 Empirical Preprint)}
\begin{document}
\maketitle
\begin{abstract}
We reconstruct the Triassic-to-Anthropocene phylogenomic radiation, endogenous lignocellulase/eusocial caste/terpenoid defense enzyme evolution, and multi-compartment demographic history of the global radiation of Isoptera and sister wood-roach outgroups (spanning Basal Relictual Families [Mastotermes darwiniensis, Zootermopsis nevadensis, Z. angusticollis, Porotermes adamsoni, Hodotermes mossambicus], Drywood Kalotermitidae [Cryptotermes havilandi, C. brevis, Incisitermes minor, Neotermes koshunensis], Subterranean Rhinotermitidae [Coptotermes formosanus, C. gestroi, Reticulitermes flavipes, R. virginicus, R. lucifugus, Heterotermes validus, Schedorhinotermes breinli], and Higher Termitidae [fungus-growing Macrotermes natalensis, M. barneyi, Odontotermes obesus, humivorous Pericapritermes nitobei, Termes fatalis, Neocapritermes angusticeps, Anoplotermes parvus, and diterpene-spraying nasutes Nasutitermes corniger, Hospitalitermes medioflavus, Constrictotermes cyphergaster, Labiotermes labralis], rooted on Cryptocercus relictus / C. kyebangensis and Blattodea outgroups Periplaneta americana and Blattella germanica) across 544 active verified NCBI sequence records (plus 317 autonomously quarantined GenBank UNVERIFIED anomalies) comprising a 64-taxon, 18-locus supermatrix (L = 34,650 aligned columns), 48 nuclear functional coding sequences, and 12 regional and multi-locus cohorts (12 x N = 36 = 432 isolates). Executed end-to-end via autonomous Native Rust bio-suite-supervisor optimization across 32 candidate RuNIE chains (total runtime 753.88 s; B = 1,000 UFBoot; 25,000 MCMC generations with 20\% burn-in), RY3 purine/pyrimidine recoding under GTR+F+R4 reduces Blattodea/Isoptera compositional heterogeneity by 90.2\% (RCFV = 0.0032 vs. 0.0326, ln L = -421,583.88), dating Root Blattodea/Isoptera to 238.4 Ma, Crown Isoptera to 148.8 Ma, Crown Neoisoptera to 98.6 Ma, and Crown Higher Termitidae to 54.2 Ma. Native Rust quartet multispecies coalescent inference (runie-quartet-coalescent across 18 gene trees and 11,436,768 quartets in 45.39 ms) recovers 8,910,008 satisfied quartets (normalized Q = 0.7791, mean LPP = 0.8386) while ABBA-BABA tests quantify secondary-contact allele sharing between invasive Coptotermes formosanus and C. gestroi. Sixty-one-state Goldman-Yang codon models identify strong purifying constraint on endogenous GH9 endoglucanase/laccase cores (omega\_0 = 0.167) and juvenile hormone/vitellogenin caste regulators alongside 4D Bayes Empirical Bayes diversifying selection on Toll/mound-hypoxia receptors (peak BEB omega = 7.12) and CYP4G15/CYP6K1 hydrocarbon/terpenoid monooxygenase pockets (foreground omega\_fg = 1.500, peak BEB omega = 2.75). Finally, 12 coupled Bayesian coalescent skylines contrast Pleistocene savanna/forest refugia and 19th-20th century maritime timber bridgehead bottlenecks against sharp locus-specific Ne compressions at nuclear GH9 cellulase and hexamerin caste loci relative to neutral mitochondrial COX1/COII (1Ne) and Katoh E-INS-i-aligned ribosomal 16S/28S/ITS rDNA (4Ne) reservoirs.
\end{abstract}
\section{Introduction}
We reconstruct the Triassic-to-Anthropocene phylogenomic radiation, endogenous lignocellulase and eusocial caste evolution, and demographic history of the global radiation of Isoptera and sister wood-roach outgroups across 544 active verified NCBI sequence records (64 stratified Blattodea and Isoptera 16-locus species OTUs assembled via \texttt{bio-suite-seqnet}, 48 nuclear functional CDS orthologs across 4 physiological gene families, and 12 regional and multi-locus cohorts totaling 432 isolates; $L = 34,650$ aligned supermatrix columns) plus 317 autonomously quarantined sequence slices logged in the persistent SQLite \texttt{quarantine\_ledger} using 100\% Native Rust tools in \texttt{bio-suite-supervisor} \cite{ref1,ref2,ref3}.
\section{Autonomous Pipeline Supervision and 32-Chain RuNIE Sweep}
Across 32 candidate Native Rust chains spanning four data-hygiene and recoding regimes, purine/pyrimidine \texttt{RY3} recoding under \texttt{GTR+F+R4} achieved $\ln L = -421,583.88$ and reduced compositional heterogeneity to $\mathrm{RCFV} = 0.0032$ in 753.88 s end-to-end execution.
\section{Calibrated Phylogenomics, Quartet Coalescent and Demographic Skylines}
Fossil-calibrated \texttt{TimeTreeLSD2} dating places Root Blattodea/Isoptera at 238.40 Ma, Crown Isoptera at 148.80 Ma, Crown Neoisoptera at 98.60 Ma, and Crown Termitidae at 54.20 Ma \cite{ref2,ref3,ref14}. Native 18-locus Quartet Multispecies Coalescent inference across 11,436,768 quartets recovers normalized $Q = 0.7791$ (mean LPP = 0.8386) \cite{ref15}, while 61-state Goldman--Yang codon models identify diversifying selection in \texttt{CYP4G15} ($\omega_2 = 2.75$) \cite{ref6,ref17} and 12 coupled 25,000-generation Bayesian skylines quantify Pleistocene refugia and maritime bridgehead bottlenecks vs. multi-locus coalescent sweeps \cite{ref11,ref16}.
\begin{thebibliography}{99}
\bibitem{ref1} White, B. P. (2026). RuNIE (Rust Native Inference Engine): An Open-Source Architecture for Autonomous Bioinformatics Pipeline Synthesis, Blinded Validation, and Bayesian Scope-Depth Auditing. CRV Science Labs Preprint / Zenodo. DOI: 10.5281/zenodo.23969275 \url{https://doi.org/10.5281/zenodo.23969275}

\bibitem{ref2} Bourguignon, T., Lo, N., Cameron, S. L.,  obotn k, J., Hayashi, Y., Shigenobu, S., Watanabe, D., Roisin, Y., Miura, T., \& Evans, T. A. (2015). The evolutionary history of termites as inferred from 66 mitochondrial genomes. Molecular Biology and Evolution, 32(2), 406--421. DOI: 10.1093/molbev/msu308 \url{https://doi.org/10.1093/molbev/msu308}

\bibitem{ref3} Bucek, A.,  obotn k, J., He, S., Bezdecka, M., K   kov , B., Sillam-Duss s, D., Hyli , M., Roisin, Y., Evans, T. A., \& Bourguignon, T. (2019). Evolution of termite symbiosis informed by transcriptome-based phylogenies. Current Biology, 29(21), 3728--3734. DOI: 10.1016/j.cub.2019.08.076 \url{https://doi.org/10.1016/j.cub.2019.08.076}

\bibitem{ref4} Bourguignon, T., Lo, N., Sobotnik, J., Ho, S. Y. W., Iqbal, N., Coissac, E., Lee, M., Jendryka, M. M., Sillam-Dusses, D., Krizkova, B., Roisin, Y., \& Evans, T. A. (2017). Mitochondrial phylogenomics resolves the global spread of higher termites, ecosystem engineers of the tropics. Molecular Biology and Evolution, 34(3), 589--597. DOI: 10.1093/molbev/msw253 \url{https://doi.org/10.1093/molbev/msw253}

\bibitem{ref5} Inward, D., Beccaloni, G., \& Eggleton, P. (2007). Death of an order: a comprehensive molecular phylogenetic study confirms that termites are eusocial cockroaches. Biology Letters, 3(3), 331--335. DOI: 10.1098/rsbl.2007.0102 \url{https://doi.org/10.1098/rsbl.2007.0102}

\bibitem{ref6} Terrapon, N., Li, C., Robertson, H. M., Ji, L., Meng, X., Booth, W., Chen, Z., Childers, C. P., Glastad, K. M., Gokhale, K., Gowda, S., Hughes, D. S. T., Jones, J. W., Korb, J., Bornberg-Bauer, E., Wang, J., \& Liebig, J. (2014). Molecular traces of alternative social organization in a termite genome. Nature Communications, 5, 3636. DOI: 10.1038/ncomms4636 \url{https://doi.org/10.1038/ncomms4636}

\bibitem{ref7} Harrison, M. C., Jongepier, E., Robertson, H. M., Arning, N., Bitard-Feildel, T., Chao, H., Childers, C. P., Dinh, H., Doddapaneni, H., Dugan, S., Gowda, S., Han, Y., Hughes, D. S. T., Liebig, J., Korb, J., \& Bornberg-Bauer, E. (2018). Hemimetabolous genomes reveal molecular basis of termite eusociality. Nature Ecology \& Evolution, 2(3), 557--566. DOI: 10.1038/s41559-017-0459-1 \url{https://doi.org/10.1038/s41559-017-0459-1}

\bibitem{ref8} Poulsen, M., Hu, H., Li, C., Chen, Z., Xu, L., Otani, S., Nygaard, S., Nobre, T., Klaubauf, S., Schindler, P. M., Hauser, F., Pan, H., Yang, Z., Sonnenberg, A. S. M., de Beer, Z. W., Zhang, Y., Wingfield, M. J., Grimmelikhuijzen, C. J. P., de Vries, R. P., Korb, J., Aanen, D. K., Wang, J., Boomsma, J. J., \& Zhang, G. (2014). Complementary symbiont contributions to plant decomposition in a fungus-farming termite. Proceedings of the National Academy of Sciences, 111(40), 14500--14505. DOI: 10.1073/pnas.1319718111 \url{https://doi.org/10.1073/pnas.1319718111}

\bibitem{ref9} Watanabe, H., Noda, H., Tokuda, G., \& Lo, N. (1998). A cellulase gene of termite origin. Nature, 394(6691), 330--331. DOI: 10.1038/28527 \url{https://doi.org/10.1038/28527}

\bibitem{ref10} Chouvenc, T.,  obotn k, J., Engel, M. S., \& Bourguignon, T. (2021). Termite evolution: mutualistic associations, key innovations, and the rise of Termitidae. Cellular and Molecular Life Sciences, 78(6), 2749--2769. DOI: 10.1007/s00018-020-03728-z \url{https://doi.org/10.1007/s00018-020-03728-z}

\bibitem{ref11} Blumenfeld, A. J., Eyer, P. A., Husseneder, C., Mo, J., Johnson, L. N. L., Wang, C., Grace, J. K., Chouvenc, T., Wang, S., \& Vargo, E. L. (2021). Bridgehead effect and multiple introductions shape the global invasion history of aiette subterranean termite Coptotermes formosanus. Communications Biology, 4(1), 196. DOI: 10.1038/s42003-021-01725-x \url{https://doi.org/10.1038/s42003-021-01725-x}

\bibitem{ref12} Chouvenc, T., Li, H. F., Austin, J., Bordereau, C., Bourguignon, T., Cameron, S. L., Cancello, E. M., Constantino, R., Costa-Leonardo, A. M., Eggleton, P., Evans, T. A., Forschler, B., Grace, J. K., Husseneder, C., K e ek, J., Lee, C. Y., Lefebvre, T., Lo, N., Messenger, M., Miura, T., Neoh, K. B., Roisin, Y., Scheffrahn, R. H., Sillam-Duss s, D.,  obotn k, J., Szalanski, A., Takematsu, Y., Vargo, E. L., Yamada, A., Yoshimura, T., \& Su, N. Y. (2016). Revisiting Coptotermes (Isoptera: Rhinotermitidae): a global taxonomic road map for speciesenez of economic importance. Systematic Entomology, 41(2), 299--306. DOI: 10.1111/syen.12157 \url{https://doi.org/10.1111/syen.12157}

\bibitem{ref13} Dedeine, F., Dupont, S., Guyot, S., Matsuura, K., Wang, C., Habibpour, B., Bagn res, A. G., Mantovani, B., \& Luchetti, A. (2016). Historical biogeography of Reticulitermes termites (Isoptera: Rhinotermitidae) inferred from analyses of mitochondrial and nuclear loci. Molecular Phylogenetics and Evolution, 94, 778--790. DOI: 10.1016/j.ympev.2015.10.020 \url{https://doi.org/10.1016/j.ympev.2015.10.020}

\bibitem{ref14} Engel, M. S., Grimaldi, D. A., \& Krishna, K. (2009). Termites (Isoptera): their phylogeny, classification, and rise to ecological dominance. American Museum Novitates, 3650, 1--27. DOI: 10.1206/651.1 \url{https://doi.org/10.1206/651.1}

\bibitem{ref15} Mirarab, S., \& Warnow, T. (2015). ASTRAL-II: Coalescent-based species tree estimation with many hundreds of taxa and thousands of genes. Bioinformatics, 31(12), i44--i52. DOI: 10.1093/bioinformatics/btv234 \url{https://doi.org/10.1093/bioinformatics/btv234}

\bibitem{ref16} Drummond, A. J., Rambaut, A., Shapiro, B., \& Pybus, O. G. (2005). Bayesian coalescent inference of past population dynamics from molecular sequences. Molecular Biology and Evolution, 22(5), 1185--1192. DOI: 10.1093/molbev/msi103 \url{https://doi.org/10.1093/molbev/msi103}

\bibitem{ref17} Yang, Z., Wong, W. S. W., \& Nielsen, R. (2005). Bayes empirical Bayes inference of amino acid sites under positive selection. Molecular Biology and Evolution, 22(4), 1107--1118. DOI: 10.1093/molbev/msi097 \url{https://doi.org/10.1093/molbev/msi097}

\end{thebibliography}
\end{document}
